\section{第9章：振动（9.5练习题）}
\sourceproblem{选择1}{163}{169}
\problemtext 判断简谐振子的加速度大小和位移大小及方向关系。
\workings $a=-\omega^2x$，非零位移处加速度大小与位移大小成正比且方向相反，选B。
\crosscheck 平衡点两者均为零，无方向可言；关系式处处有效。
\sourceproblem{选择2}{163}{169}
\problemtext 水平弹簧振子运动到最大位移时，质量 $m_0$ 的泥块竖直落下粘住原质量m，判断振幅与周期变化。
\workings 在光滑水平面上，碰撞前后水平动量守恒，瞬时水平速度均为零，平衡位置不变，故新振幅仍为原最大位移A。周期从 $2\pi\sqrt{m/k}$ 变为 $2\pi\sqrt{(m+m_0)/k}$，增大，选C。
\crosscheck 碰撞不要求总机械能守恒；竖直落下的动能可耗散，不改变该水平初始条件。
\sourceproblem{选择3}{163}{169}
\problemtext 均匀轻弹簧原周期T，剪成两等长半段并联悬挂同一物体，求新旧周期比（选项续PDF170）。
\workings 每半段劲度为2k，并联等效4k，故 $T'/T=\sqrt{k/(4k)}=1/2$，选D。
\crosscheck 依赖均匀弹簧和线性小形变；两半若串联则恢复原k。
\sourceproblem{选择4}{164}{170}
\problemtext 两轻弹簧劲度 $k_1,k_2$ 并联悬挂质量m，求竖直平动固有频率。
\workings 在平衡点附近共同位移x产生合回复力 $-(k_1+k_2)x$，故 $\nu=(1/2\pi)\sqrt{(k_1+k_2)/m}$，选A。
\crosscheck 假定物体不转动、两弹簧平行且线性；重力改变平衡位置而不改变频率。
\sourceproblem{选择5}{164}{170}
\problemtext 从给定加速度位移函数中选简谐运动形式。
\workings 只有 $a=-10x$ 具有恒定负比例系数，选C，角频率为 $\sqrt{10}$ rad/s（SI）。
\crosscheck 正比例是反回复力，平方项非线性。
\sourceproblem{选择6}{164}{170}
\problemtext 图中三曲线1、2、3分别代表简谐运动位移、速度、加速度，辨别对应关系。
\workings 曲线1和3始终反相，曲线1在初始过零处斜率为负，曲线2此刻为负极值；所以1为x、2为v、3为a，选D。
\crosscheck 在曲线1的负极值时曲线2过零且曲线3为正极值，符合 $v=\dot x,a=-\omega^2x$；不能仅按纵向振幅判断物理量。
\sourceproblem{选择7}{164}{170}
\problemtext $x=A\cos(\omega t+\pi/4)$，求半周期时加速度。
\workings $t=T/2$ 时相位 $5\pi/4$，$a=-A\omega^2\cos(5\pi/4)=A\omega^2/\sqrt2$，选B。
\crosscheck 此时位移为负，回复加速度为正。
\sourceproblem{选择8}{164}{170}
\problemtext $x=A\cos(\omega t+\phi)$，求半周期时速度。
\workings $v=-A\omega\sin(\pi+\phi)=A\omega\sin\phi$，选B。
\crosscheck 半周期后速度反向，大小与初速度相同。
\sourceproblem{选择9}{164}{170}
\problemtext 质量 $m_1,m_2$ 的球连接同一轻弹簧，固定 $m_1$ 时 $m_2$ 振动频率 $\nu_2$，改固定 $m_2$ 时求 $\nu_1$。
\workings 同一劲度k给 $\nu_i=\sqrt{k/m_i}/(2\pi)$，所以 $\nu_1=\nu_2\sqrt{m_2/m_1}$，选D。
\crosscheck 不是两球均自由振动时的约化质量模型。
\sourceproblem{选择10}{164}{170}
\problemtext 水平弹簧振子拉离平衡位置5 cm后静止释放，开始向x负向运动，求余弦式振幅和初相。
\workings 初速度为零，因释放后向负方向运动，初始回复加速度为负，故 $x_0=+5$ cm，$A=5$ cm、$\phi=0$（模 $2\pi$），选A。
\crosscheck 若初始位移为负极值，则释放后应向正方向运动。
\sourceproblem{选择11}{164}{170}
\problemtext 简谐振子初位移 $A/2$，向x正向运动，求余弦式初相。
\workings $\cos\phi=1/2$ 且 $v_0=-A\omega\sin\phi>0$，故 $\phi=-\pi/3$（模 $2\pi$），选A。
\crosscheck $+\pi/3$ 虽满足初位移却对应负速度。

\sourceproblem{选择12}{165}{171}
\problemtext 振幅A、周期T的振子初位移 $A/2$，向负x运动，选振动方程。
\workings $\omega=2\pi/T$，$\cos\phi=1/2$，且 $v_0=-A\omega\sin\phi<0$，故 $\phi=\pi/3$，$x=A\cos(2\pi t/T+\pi/3)$，选D。
\crosscheck 负初相会给出相反速度；原页分数为A/2。
\sourceproblem{选择13}{165}{171}
\problemtext 两等频等幅振子反向相遇时离平衡位置均为半振幅，求相位差。
\workings 在正半振幅处两相位为 $\pm\pi/3$，负半振幅处为 $\pm2\pi/3$。两种相遇给同一无向相位差 $2\pi/3$，选D。
\crosscheck 对非同相振动，两次相邻相遇相隔半周期，位移互为相反数；若把原题两次位移均为正A/2按字面理解，则不能是相邻两次相遇。按位移大小理解；有向相位差可为正负 $2\pi/3$。
\sourceproblem{选择14}{165}{171}
\problemtext 两等频等幅振子，第一过零向负，第二处于 $-A/2$ 且向负，求第二相对第一相位差。
\workings 同时刻相位分别为 $\pi/2,2\pi/3$，差为 $\pi/6$，选C。
\crosscheck 两速度均负，不能选择第二相位 $4\pi/3$。
\sourceproblem{选择15}{165}{171}
\problemtext A从正位移回到零时B在正最大位移，比较两同频振子的相位。
\workings 此时 $\theta_A=\pi/2,\theta_B=0$（模 $2\pi$），A超前B $\pi/2$，选B。
\crosscheck A此时速度为负，故相位为正半周内的过零相位。
\sourceproblem{选择16}{165}{171}
\problemtext 周期T、振幅A的振子由零点至离零点 $A/2$ 的最短时间。
\workings 以正向过零为初态，$x=A\sin\omega t$，首次达到半幅时 $\omega t=\pi/6$，故 $\Delta t=T/12$，选D。
\crosscheck 不能将简谐运动视为匀速运动；向负侧结果相同。
\sourceproblem{选择17}{165}{171}
\problemtext $x=0.04\cos(2\pi t+\pi/3)$ m，求初时至 $x=-0.02$ m且速度为正的最短时间。
\workings 目标相位为 $4\pi/3$，初相 $\pi/3$，最小正相位增量 $\pi$，故 $\Delta t=1/2$ s，选D。
\crosscheck 相位 $2\pi/3$ 虽同位移但速度为负。
\sourceproblem{选择18}{165}{171}
\problemtext 弹簧振子质量m、劲度k、振幅A，初时正向过零，选振动方程。
\workings $\omega=\sqrt{k/m}$，初相 $-\pi/2$，故 $x=A\cos(\sqrt{k/m}\,t-\pi/2)$，选B。
\crosscheck $v(0)=A\sqrt{k/m}>0$。
\sourceproblem{选择19}{165}{171}
\problemtext 向下为正的竖直弹簧盘在最低点粘上物体，速度不突变，新平衡点下移量d小于原振幅A，求新初相范围（选项续PDF172）。
\workings 碰撞处原盘速度为零，新位移 $x'_0=A-d>0$。此为新振动的正极值，$A'=A-d$，故 $\phi'=0$（模 $2\pi$）。
\crosscheck 原题各选项为相位区间，严格结果在A选项的零端点；若区间均为开区间则没有正确选项，若把零包含在A内则选A。不可由位置为正擅自推为严格第一象限。
\sourceproblem{选择20}{166}{172}
\problemtext 单摆向正位移拉开小角度后由静止释放，求余弦式初相。
\workings 初始正极值且速度零，$\phi=0$，选C。
\crosscheck 采用小角近似和释放时刻为零。
\sourceproblem{选择21}{166}{172}
\problemtext 第一振子 $x_1=A\cos(\omega t+\alpha)$ 从正位移回零时，第二等幅等周期振子在正极值，求第二振动式。
\workings 此时第一相位 $\pi/2$，第二相位0，故 $x_2=A\cos(\omega t+\alpha-\pi/2)$，选B。
\crosscheck 选项C的 $-3\pi/2$ 与 $+\pi/2$ 等价，却使第二处在负极值。
\sourceproblem{选择22}{166}{172}
\problemtext 长l均匀杆以端点水平轴为支点作小幅复摆运动，求周期。
\workings $J=ml^2/3$，质心距轴 $l/2$，$J\ddot\theta=-mg(l/2)\sin\theta\simeq-mg(l/2)\theta$。故 $T=2\pi\sqrt{2l/(3g)}$，选C。
\crosscheck 有效摆长为 $J/[m(l/2)]=2l/3$，不是杆长。
\sourceproblem{选择23}{166}{172}
\problemtext 弹簧振子和小角单摆从地面移至月面，比较固有周期。
\workings 理想弹簧振子 $T_1=2\pi\sqrt{m/k}$ 不依g；单摆 $T_2=2\pi\sqrt{l/g}$，月面g较小故 $T'_2>T_2$，选D。
\crosscheck 假定质量、劲度和摆长不变。
\sourceproblem{选择24}{166}{172}
\problemtext 简谐振动频率为 $\nu$，求动能变化频率。
\workings $K=mA^2\omega^2[1-\cos(2\omega t+2\phi)]/4$，非恒定部分频率为 $2\nu$，选B。
\crosscheck 半个振动周期动能即重复。
\sourceproblem{选择25}{166}{172}
\problemtext $x=5\cos(2\pi t+\pi/2)$ m，判断初始运动状态。
\workings $x_0=0,v_0=-10\pi$ m/s，$a_0=-\omega^2x_0=0$，初相 $\pi/2$，总能量非零，选C。
\crosscheck 零位移处动能最大而非能量为零。
\sourceproblem{选择26}{166}{172}
\problemtext 两同劲度弹簧挂质量 $m_1>m_2$，等振幅，比较周期和振动能量。
\workings $T=2\pi\sqrt{m/k}$，故第一周期更长；以各自平衡点为零的振动能 $E=kA^2/2$ 相同，选C。
\crosscheck 不能以质量较大推断等幅振动能较大，角频率亦改变。
\sourceproblem{选择27}{166}{172}
\problemtext 振幅A的弹簧振子动能等于振动势能时的位置（选项续PDF173）。
\workings 以平衡点为势能零点，$K=k(A^2-x^2)/2$，$U=kx^2/2$。相等给 $x=\pm A/\sqrt2$，选C。
\crosscheck 两点能量均为总振动能的一半。
\sourceproblem{选择28}{167}{173}
\problemtext 两同向同频等幅A振动合幅亦为A，求初相差。
\workings $A^2=2A^2(1+\cos\Delta\phi)$，故 $\cos\Delta\phi=-1/2$，无向相位差 $2\pi/3$，选C。
\crosscheck 有向相位差为 $\pm2\pi/3$ 模 $2\pi$。
\sourceproblem{选择29}{167}{173}
\problemtext 两同向振动 $x_1=A_1\cos\omega t,x_2=A_2\sin\omega t$，求合幅。
\workings 两相量正交，合幅 $\sqrt{A_1^2+A_2^2}$，选C；合初相为 $-\arctan(A_2/A_1)$（正振幅约定）。
\crosscheck 展开合余弦恢复正的正弦项。
\sourceproblem{选择30}{167}{173}
\problemtext 两同向振动幅值均0.04 m、角频率6 rad/s，初相分别 $\pm\pi/3$，求合振动。
\workings 余弦和差公式给 $x=2(0.04)\cos(\pi/3)\cos6t=0.04\cos6t$ m，选B。
\crosscheck 正弦分量抵消，合幅并非两振幅直接相加。
\sourceproblem{填空1}{167}{173}
\problemtext 竖直弹簧挂质量m时周期2 s，增加2 kg后周期3 s，求m、k。
\workings $T^2=4\pi^2m/k$，故 $(m+2)/m=9/4$，$m=1.6$ kg，$k=\pi^2m=1.6\pi^2$ N/m（约15.79 N/m）。
\crosscheck 代回两个质量分别得2 s、3 s。
\sourceproblem{填空2}{167}{173}
\problemtext 振幅0.15 m、角频率 $6\pi$ rad/s、初相 $\pi/2$ 的余弦振动。
\workings $x=0.15\cos(6\pi t+\pi/2)$ m。
\crosscheck 初位移零，初速度为 $-0.9\pi$ m/s。
\sourceproblem{填空3}{167}{173}
\problemtext 竖直弹簧劲度k挂质量m，在自然长度处静止释放，求频率和振幅。
\workings 平衡伸长 $mg/k$，相对平衡点初位移绝对值等于此伸长，故 $A=mg/k$，$\nu=\sqrt{k/m}/(2\pi)$。
\crosscheck 若向下为正，初位移为 $-mg/k$；最大伸长为 $2mg/k$。
\sourceproblem{填空4}{167}{173}
\problemtext 质量m、劲度k的理想弹簧振子固有周期。
\workings $m\ddot x=-kx$，$T=2\pi\sqrt{m/k}$。
\crosscheck k的单位为N/m，根号内量纲为时间平方。
\sourceproblem{填空5}{167}{173}
\problemtext 根据三条标为 $x_1,x_2,x_3$ 的振动曲线写方程。
\workings 高清图显示三者幅值10 cm，峰值间隔2 s，故 $\omega=\pi$ rad/s。$x_1$ 在初时正极值，$x_2$ 初时正向过零，$x_3$ 初时负极值；按厘米为单位，
\[
x_1=10\cos\pi t,\quad x_2=10\cos(\pi t-\pi/2),\quad x_3=10\cos(\pi t+\pi).
\]
\crosscheck 三条曲线首次正峰分别在0、0.5、1 s，随后相隔2 s重复；图中 $x_1$ 标签位于第二正峰处。
\sourceproblem{填空6}{167}{173}
\problemtext 弹簧挂M时t秒振动n次，再挂质量m，求新周期。
\workings 原周期 $T=t/n$，同一劲度下 $T'/T=\sqrt{(M+m)/M}$，所以 $T'=(t/n)\sqrt{1+m/M}$。
\crosscheck $m=0$ 时恢复原周期，假设弹簧保持线性。
\sourceproblem{填空7}{167}{173}
\problemtext 同劲度弹簧的两振子质量比4:1，求周期比。
\workings 周期与质量平方根成正比，故 $T_1:T_2=2:1$。
\crosscheck 频率比反为1:2。
\sourceproblem{填空8}{167}{173}
\problemtext 0.4 N力使轻弹簧伸长0.20 m，求周期 $2\pi$ s所需悬挂质量。
\workings $k=F/\Delta l=2$ N/m，$m=k(T/2\pi)^2=2$ kg。
\crosscheck 假定加载后弹簧仍服从胡克定律。
\sourceproblem{填空9}{167}{173}
\problemtext 两振动曲线如图，求 $x_1,x_2$ 的最大速率之比。
\workings 高清图读得 $A_1=1$ cm、$A_2=2$ cm；$T_1=2$ s、$T_2=4$ s。$v_{\max}=2\pi A/T$，故 $v_{1,\max}:v_{2,\max}=1:1$。
\crosscheck 两者最大速率均为 $\pi$ cm/s；不单凭振幅比判断。
\sourceproblem{填空10}{168}{174}
\problemtext 图示水平振子在负向过零及正最大位移时对应的曲线标记。
\workings 负向过零满足 $x=0,v=-\omega A,a=F=0$，对应b、f。正极值满足 $x=A,v=0,a=-kA/m,F=-kA$，对应a、e。
\crosscheck c虽为极值但加速度和力为正；O、d为正向过零。
\sourceproblem{填空11}{168}{174}
\problemtext 无阻尼自由简谐振动的周期、频率及振幅、初相分别由何决定。
\workings 周期频率由系统固有参数（如m、k）决定；对给定系统，振幅与初相由初始位移和初始速度决定。
\crosscheck 零振幅时初相无物理定义；非线性系统频率可能依幅值。
\sourceproblem{填空12}{168}{174}
\problemtext 振幅A、周期T，初位移 $A/2$ 且负向运动，求初相和方程。
\workings $\cos\phi=1/2,\sin\phi>0$，故 $\phi=\pi/3$，$x=A\cos(2\pi t/T+\pi/3)$。
\crosscheck 与本章选择12同一初态，独立保留原书题号。
\sourceproblem{填空13}{168}{174}
\problemtext 沿y轴振动，初位移 $-A/2$ 且正向运动，求初相和方程。
\workings $\cos\phi=-1/2,\sin\phi<0$，故 $\phi=-2\pi/3$（模 $2\pi$），$y=A\cos(2\pi t/T-2\pi/3)$。
\crosscheck 初速度为 $\sqrt3 A\omega/2>0$。
\sourceproblem{填空14}{168}{174}
\problemtext 初位移 $A/2$ 且正向运动，4 s后首次通过平衡点，求初相及角频率。
\workings $\phi=-\pi/3$；先经过正峰再负向过零，目标相位 $\pi/2$，故 $4\omega=5\pi/6$，$\omega=5\pi/24$ rad/s。
\crosscheck 首次过零不能取初相之前的 $-\pi/2$；周期为9.6 s。
\sourceproblem{填空15}{168}{174}
\problemtext 长l的单摆，悬点正下方 $l/2$ 处有钉，小幅摆动求完整周期。
\workings 摆线在一侧自由长l，另一侧绕钉的有效摆长 $l/2$。两侧各耗相应单摆的半周期，故 $T=\pi\sqrt{l/g}+\pi\sqrt{l/(2g)}$。
\crosscheck 钉在最低点处改变支点，理想无摩擦接触且切换时速度连续；不能平均摆长后代入单摆公式。
\sourceproblem{填空16}{168}{174}
\problemtext 图示单摆初态(a)正向过零、(b)负向过零、(c)左侧静止极值，角坐标右为正，写三个余弦方程。
\workings 设A为角振幅，$\omega=2\pi/T$，则 $\theta_a=A\cos(\omega t-\pi/2)$，$\theta_b=A\cos(\omega t+\pi/2)$，$\theta_c=A\cos(\omega t+\pi)$。
\crosscheck 图c位于负位移侧，初速度为零；若A定义为弧长振幅，角振幅须改为A/l。
\sourceproblem{填空17}{168}{174}
\problemtext 理想弹簧振子振幅增至两倍，求振动能与频率的倍数。
\workings $E=kA^2/2$，故能量为4倍；$\nu=\sqrt{k/m}/(2\pi)$ 不变，为1倍。
\crosscheck 假定质量、劲度不变且振幅仍在线性范围。
\sourceproblem{填空18}{168}{174}
\problemtext $x_1=0.06\cos(5t+\pi/2)$、$x_2=0.02\sin(\pi-5t)$（SI），求合幅与初相。
\workings $x_1=-0.06\sin5t$，$x_2=0.02\sin5t$，合振动 $x=-0.04\sin5t=0.04\cos(5t+\pi/2)$ m，故 $A=0.04$ m，$\phi=\pi/2$。
\crosscheck 两振动反相，较大分振动决定合相位；初速度为 $-0.20$ m/s。
